\documentclass[10pt,a4paper]{amsart}
\usepackage[margin=19mm]{geometry}
\usepackage{amsmath,amssymb,mathtools}
\usepackage[scr=rsfs,cal=euler]{mathalfa}
\usepackage{xfrac}
\usepackage[unicode,hidelinks,bookmarks=false]{hyperref}
\newcommand{\ie}{i.e.\ }
\newcommand{\cf}{cf.\ }
\newcommand{\resp}{respectively\ }
\newcommand{\Subsection}{\subsection}
\providecommand{\nobreakdash}{\nobreak\hbox{-}}
\setlength{\emergencystretch}{2em}
\multlinegap0pt
\usepackage{enumitem}
\newtheorem{theorem}{Theorem}
\newtheorem{proposition}{Proposition}
\newtheorem{lemma}{Lemma}
\newtheorem{corollary}{Corollary}
\newtheorem{claim}{Claim}
\newcommand{\N}{\mathbb{N}}
\newcommand{\Q}{\mathbb{Q}}
\newcommand{\abs}[1]{\left\lvert#1\right\rvert}
\DeclareMathOperator{\osc}{osc}
\DeclareMathOperator{\diam}{diam}
\DeclareMathOperator{\dist}{dist}
\DeclareMathOperator{\lgth}{lh}
\title[main theorem]{Continuous functions on limits of F-decomposable systems: an LUR renorming of $C(X)$ for $F_d$-compacta (main theorem)}
\author{Todor Manev}
\date{Source: arXiv:2503.16626v3 (CC BY 4.0)}
\begin{document}
\maketitle

\noindent\emph{Source and licence.} Continuous functions on limits of F-decomposable systems, by Todor Manev. Version arXiv:2503.16626v3
(CC BY 4.0), distributed by arXiv under the Creative Commons Attribution 4.0 International licence,
http://creativecommons.org/licenses/by/4.0/, as recorded in the arXiv licence metadata for this exact
version and shown on the abstract page https://arxiv.org/abs/2503.16626v3. Full licence notice:
\texttt{licenses/arxiv-2503.16626-CC-BY-4.0.txt}.\par\smallskip

\noindent\textit{Editorial scope.} Counted unit: the article's main theorem (source0.tex lines 382--384): if $X$ is the limit of an F-decomposable inverse system and an $\mathrm{F}_d$-compact of spectral height $\gamma<\omega_1$, then $C(X)$ admits an equivalent $\tau_p$-lower semicontinuous LUR norm. The article's complete original proof of it is delivered with it (source0.tex lines 386--457, one \texttt{proof} environment of 72 lines, adjacent to the statement, with no gap and no deferral). No mathematics is written, repaired or re-derived: the statement and the proof are the article's own lines, in the article's own notation, and no retained line is substituted, re-flowed or omitted, so nothing is removed and the package carries no omission record. Retained uncounted as the source-local context the proof needs: the metrizability criterion for fully closed mappings (lines 185--192, source label \texttt{Fed\_metr}), the characterization of LUR renormings by countable decompositions (lines 202--207, source label \texttt{LUR\_char}), the metric projection proposition (lines 123--125, source label \texttt{metr\_proj}), the finite-branching lemma (lines 347--349, source label \texttt{finite\_branching}) and the section heading under which the construction sits (line 236). Every cross-reference of every retained line resolves inside this document.

Bibliography: the retained lines cite four works whose entries are transcribed verbatim from the article's own bibliography block in source0.tex (the entries numbered \texttt{fedFC}, \texttt{holmes\_opt}, \texttt{mot97} and \texttt{motv}) and delivered as a four-entry bibliography with short editorial labels; the mapping is recorded in the item's reference mapping. No other bibliography entry is transcribed and no retained line cites anything else.

Editorial formatting only, in addition: an AMS \texttt{amsart} preamble that restates the article's own declarations and macros used by the retained lines -- the environments \texttt{theorem}, \texttt{proposition}, \texttt{lemma} and \texttt{corollary}, each on its own counter so that cross-references name the right kind of statement, and the macros \texttt{\textbackslash N}, \texttt{\textbackslash Q}, \texttt{\textbackslash abs}, \texttt{\textbackslash osc}, \texttt{\textbackslash diam}, \texttt{\textbackslash dist} and \texttt{\textbackslash lgth} copied or adapted from the article's own preamble (source0.tex lines 39, 41, 46, 47, 49 and 50, with \texttt{\textbackslash abs} taking the definition the article's own \texttt{commath} package gives it). The retained environments print in this document as Proposition 1, Proposition 2, Theorem 3, Lemma 4 and Theorem 5 in order of appearance, whereas the article prints the same five items with the numbering of its own full document; the numbered displays of the retained lines print consecutively in order of appearance, whereas the article numbers its equations per section.

The complete native source of the article as distributed in the arXiv e-print for this exact version is bundled unchanged as \texttt{sources/175132/source0.tex}: it is byte identical to the e-print file \texttt{f-decomposable.tex} except that the pipeline's mechanical concatenation prepends one header comment line naming it. The article was typeset with the standard \texttt{article} class; no publisher class or style file is shipped in the e-print and none is redistributed or used.
\begin{proposition}[see e.g. {\cite[p. 173]{holmes_opt}}] \label{metr_proj}
	Let $X$ and $Y$ be Hausdorff compact spaces and $\varphi: X \to Y$ a continuous surjection. Let $\varphi^0$ denote the natural embedding of $C(Y)$ into $C(X)$, that is, $(\varphi^0 g)(x) := g(\varphi(x))$, for $g \in C(Y)$. Then, whenever $f \in C(X)$, there exists a function $g \in C(Y)$ such that $\|f- \varphi^0 g\| = \dist(f, \varphi^0 C(Y)) = \frac{1}{2} \sup \{\osc_{\varphi^{-1}(y)} f: y \in Y\}$.
\end{proposition}

\begin{proposition} [{\cite[II.3.10]{fedFC}}] \label{Fed_metr}
	Let $f: X \to Y$ be a fully closed mapping between Hausdorff compacta. Then $X$ is metrizable if and only if the following conditions hold:
	\begin{itemize}[noitemsep]
		\item $Y$ is metrizable;
		\item All the fibers $f^{-1}(y)$ are metrizable;
		\item The set of nontrivial fibers $\left\{y \in Y: \left|f^{-1}(y)\right| \geq 2\right\}$ is countable.
	\end{itemize}
\end{proposition}

\begin{theorem}[\cite{mot97}]\label{LUR_char}
	Let $E$ be a Banach space and $F$ a norming subspace of its dual. Then $E$ admits an equivalent $\sigma(E,F)$-lower semicontinuous LUR norm if and only if for any $\epsilon >0$ there exists a countable decomposition $E = \bigcup_{n \in \N} E_n$ such that for all $n \in \N$ and $x \in E_n$ there exists a $\sigma(E,F)$-open halfspace $H$ containing $x$ and satisfying:
	\begin{equation*}
		\|\cdot\| \text{-} \diam \left(E_n \cap H \right) < \epsilon.
	\end{equation*}
\end{theorem}

\subsection{Construction} \label{constr}

\begin{lemma} \label{finite_branching}
	Let $S:= \{X_\alpha, \pi_\alpha^\beta, \alpha, \beta \in \gamma\}$ be an F-decomposable system and $X$ be its limit. Then for any $f \in C(X)$,  $\nu_f$ is finite.
\end{lemma}

\begin{theorem}\label{main_thm}
	Let $X = \lim\limits_{\leftarrow} \{X_\alpha, \pi_\alpha^\beta, \alpha, \beta \in \gamma\}$ be an $\mathrm{F}_d$-compact of spectral height $\gamma < \omega_1$. Then $C(X)$ admits an equivalent $\tau_p$-lsc LUR norm.
\end{theorem} 

\begin{proof}
	Let $\epsilon > 0$. We will find a countable decomposition of the space $C(X)$ that satisfies the conditions of Theorem \ref{LUR_char}. We shall proceed in several steps.
	
	As for any $f \in C(K)$ we have $\pi^\alpha_\beta\left(M^{f,\epsilon}_\alpha\right) \subset M^{f,\epsilon}_\beta$, we can once again apply the construction from Section \ref{constr}. Denote the resulting compact space $K^{f, \epsilon}$. The full closedness of the bonding mappings implies that the sets $M^{f,\epsilon}_\alpha$ are finite for all $\alpha \in \gamma$. Thus the space $K^{f,\epsilon}$ is metrizable. The above assertion allows us to start with the first decomposition as follows:
	\begin{equation} \label{dec1}
		C(X) = \bigcup_{\nu \in \gamma^{<\omega}} E_\nu: \quad f \in E_\nu, \text{ if } \nu_f = \nu.
	\end{equation}
	
	In what follows we shall denote by $\lgth(\nu)$ the length of a finite sequence. For the subsequent decomposition we make use of the fact that each $M^{f, \epsilon}_{\nu_i}$ is finite:
	\begin{equation} \label{dec2}
		E_\nu = \bigcup_{k \in \N^{\lgth(\nu)}} E_{\nu,k}: \quad f \in E_{\nu,k}, \text{ if } \abs{M^{f,\epsilon}_{\nu_i}} = k_i; \quad i = 1, \dots, \lgth(\nu). 
	\end{equation}
	
	Often the construction of a countable decomposition satisfying the hypotheses of Theorem \ref{LUR_char} requires passing from a weak (pointwise) neighborhood of a given point to an open (pointwise open) halfspace. In this note we will do this with a technique similar to the one in \cite[Lemma 4.22]{motv}. In particular, we have the following:
	\begin{equation}
		\begin{aligned}\label{dec3}
			&E_{\nu, k} = \bigcup_{r \in \N} E_{\nu, k, r}: \quad f \in E_{\nu, k, r}, \text{ if }\\
			&\epsilon - \max\left\{\osc_{\pi_{\nu_i}^{-1}(x)} f: \;  i \in \{1, \dots, \lgth(\nu)\}, x \in \left(X_{\nu_i} \setminus M^{f, \epsilon}_{\nu_i}\right)\right\}> \frac{1}{r}.
		\end{aligned}
	\end{equation}
	
	Note that the maximum above is well defined by the full closedness of $\pi_{\nu_i}$. The next decomposition will be used bellow together with (\ref{dec3}) to combine a finite number of functionals into one. Essentially, it allows us to choose the functionals associated with a function $f \in C(X)$ in such a way that they almost attain their supremum on the corresponding piece of the decomposition at $f$. We will need to well-order the spaces $X_\alpha$ to ensure that what follows is well defined.
	\begin{equation}\label{dec4}
		E_{\nu, k, r} = \bigcup_{q \in \Q^{\sum k_i}} E_{\nu, k, r, q}: \quad \osc_{\pi_{\nu_i}^{-1}(x_j)} f \in \left(q_j - \frac{1}{3 r \sum_i k_i}, q_j + \frac{1}{3 r \sum_i k_i}\right),
	\end{equation}
	where $x_j$ are the points of $M^{f, \epsilon}_{\nu_i}$, ordered first by the index of their space $\nu_i$ and then by the well-ordering of $X_{\nu_i}$. The index $q$ is thus a multi index with $\sum_{i = 1}^{\lgth(\nu)} k_i$ coordinates.
	
	For the last decomposition we will use the construction of Section \ref{constr} with the sets $M^{f,\epsilon}_\alpha$. Denote the resulting compact space $K_f^\epsilon$. We can inductively show that this space is metrizable. Indeed, if we suppose some lower level space $K_{f, \alpha}^\epsilon$ is metrizable, the same will be true for $K_{f, \alpha +1}^\epsilon$ by Proposition \ref{Fed_metr} and the fact that the set of nontrivial fibers is finite. Now if $\alpha \in \gamma$ is a limit ordinal and $K_{f, \beta}^\epsilon$ are metrizable for all $\beta \in \alpha$, $K_{f, \alpha}^\epsilon$ will also be metrizable as it is homeomorphic to a subspace of a countable product of metrizable spaces. We thus have that $C(K_f^\epsilon)$ is separable and we can fix a countable dense subset $\{h_p\}_{p \in \N} \subset C(K_f^\epsilon)$. Then we set:
	\begin{equation}\label{dec5}
		E_{\nu, k, r, q} = \bigcup_{p \in \N} E_{\nu, k, r, q, p}: \quad  f \in E_{\nu, k, r, q,p}, \text{ if } f \in B_{\epsilon}(h_p).
	\end{equation}
	This is possible because $q^\epsilon_f: X \to K^\epsilon_f$ is a continuous surjection and by Proposition \ref{metr_proj} we have $\dist(f, C(K^\epsilon_f)) = \frac{1}{2} \sup\{\osc_{(q^\epsilon_f)^{-1}(x)}: x \in K^\epsilon_f\} < \epsilon$.
	
	We are now ready to select the functional giving the required halfspace. For $f \in E_{\nu,k,r,q,p}$ and $x \in M^{f, \epsilon}_{\nu_i}$ for some $i \leq \lgth(\nu)$, we choose $\overline{x}, \underline{x} \in \pi_{\nu_i}^{-1} (x)$ such that $f(\overline{x}) - f(\underline{x}) = \osc_{\pi_{\nu_i}^{-1}(x)} f$. We can then define the functional $\mu^f_x \in \left(C(X), \tau_p\right)^*$ as follows:
	\begin{equation*}
		\mu^f_x(g) := g(\overline{x}) - g(\underline{x}).
	\end{equation*}
	Finally, we define
	\begin{equation*}
		\mu^f := \sum \left\{\mu^f_x : \; x \in \bigcup\limits_{i = 1}^{\lgth(\nu)} M^{f, \epsilon}_{\nu_i} \right\}.
	\end{equation*}
	The required halfspace will then be given by:
	\begin{equation}
		H_f := \left\{g \in C(X): \; \mu^f (g) > \mu^f (f) - \frac{1}{3r}\right\}.
	\end{equation}
	
	\begin{claim}
		If $g \in E_{\nu,k,r,q,p} \cap H_f$, then $\mu^f_x(g) > \epsilon$ for all $x \in \bigcup_{i = 1}^{\lgth(\nu)} M^{f, \epsilon}_{\nu_i} $.
	\end{claim}
	
	\begin{proof}
		Observe first that if $\mu^f_x (g) \geq \epsilon$, for some $i \leq \lgth(\nu)$ and $x \in X_{\nu_i}$, then $x \in M^{g,\epsilon}_{\nu_i}$. As $g \in E_{\nu,k,r,q,p}$, we have that $\osc\limits_{\pi_i^{-1}(x)} g \in \left(q_j - \frac{1}{3 r \sum_i k_i}, q_j + \frac{1}{3 r \sum_i k_i}\right)$ for some $j \in \{1, \dots, \sum k_i\}$. Also, for all $j \in \{1, \dots, \sum k_i\}$, the inequality $q_j > \epsilon  - \frac{1}{3 r \sum_i k_i}$ holds.
		
		Let $\{x_1, \dots, x_l\} \subset \bigcup\limits_{i = 1}^{\lgth(\nu)} M^{f,\epsilon}_{\nu_i}$ be the subset consisting of points that satisfy $\mu^f_{x} (g) < \epsilon$, $x \in \{x_1, \dots, x_l\}$. We have:
		\begin{align*}
			&\sum \left\{\mu^f_x (g): \; x \in \left(\bigcup_{i = 1}^{\lgth(\nu)} M^{f, \epsilon}_{\nu_i}\right) \setminus \{x_1, \dots, x_l\}\right\} \\
			< \quad &\sum_{j = 1}^{\sum k_i} q_j + \left(\sum_{i = 1}^{\lgth(\nu)} k_i\right) \frac{1}{3 r \sum_{i = 1}^{\lgth(\nu)} k_i} - l\left(\min_j q_j + \frac{1}{3 r \sum_{i = 1}^{\lgth(s)} k_i}\right) \\
			< \quad& \mu^f (f) + 2 \left(\sum_{i = 1}^{\lgth(\nu)} k_i\right) \frac{1}{3 r \sum_{i = 1}^{\lgth(\nu)} k_i} - l\epsilon \\
			= \quad&\mu^f (f) + \frac{2}{3r} - l \epsilon.
		\end{align*}
		On the other hand, as $g \in H_f$, we have
		\begin{align*}
			&\sum \left\{\mu^f_x (g): \; x \in \left(\bigcup_{i = 1}^{\lgth(\nu)} M^{f, \epsilon}_{s_i}\right) \setminus \{x_1, \dots, x_l\}\right\} \\
			&> \quad \mu^f (g) - l \left(\epsilon - \frac{1}{r}\right) \\
			&> \quad \mu^f (f) - \frac{1}{3r} - l \left(\epsilon - \frac{1}{r}\right) \\
			&= \quad \mu^f (f) - \frac{1 - 3l}{3r} - l\epsilon,
		\end{align*}
		a contradiction whenever $l \geq 1$.
	\end{proof}
	
	By Lemma \ref{finite_branching} and the definition of $\nu$ the sets $M^{f, \epsilon}_\alpha$ are completely determined by the sets $M^{f, \epsilon}_{\nu_i}$. It follows that if $g \in E_{\nu,k,r,q,p} \cap H_f$, then $M^{g, \epsilon}_\alpha = M^{f, \epsilon}_\alpha$ for all $\alpha \in \gamma$ and thus $K^\epsilon_g = K^\epsilon_f$. From the decomposition (\ref{dec5}) we have that $\|f - h_p\| < \epsilon$ and $\|g - h_p\| < \epsilon$. Thus $\|f-g\| < 2 \epsilon$, which concludes the proof.
\end{proof}

\begin{thebibliography}{99}
\bibitem[Fed06]{fedFC} V. V. Fedorchuk, \emph{Fully closed mappings and their applications}, J. Math. Sci., Vol. 136, 2006, 4201--4292.
\bibitem[Hol72]{holmes_opt} R. Holmes, \emph{A Course on Optimization and Best Approximation}, Lecture Notes in Math., Vol. 257, Springer-Verlag, 1972.
\bibitem[MOT97]{mot97} A. Molt\'o, J. Orihuela, and S. Troyanski, \emph{Locally uniformly rotund renorming and fragmentability}, Proc. London Math. Soc., Vol. 75, 1997, 619--640.
\bibitem[MOTV09]{motv} A. Molt\'o, J. Orihuela, S. Troyanski, and M. Valdivia, \emph{A nonlinear transfer technique for renorming}, Lecture Notes in Math., Vol. 1951, Springer, 2009.
\end{thebibliography}
\end{document}
